\section{A15.1b2d parameter wedge: kernel triviality on
$\sigma \le \min\{d/2, R + \eta\}$}
\label{sec:p12-b2d-wedge}

The slice theorem of \S\ref{sec:p12-b2d-core-single-slice} gives
$h(x) = 0$ pointwise on $x \in (R, \min\{\sigma, d-\sigma, R+\eta\})$.
For $(R, \sigma)$ inside a specific parameter wedge, this slice
covers the whole defect strip $(R, \sigma)$, and full kernel
triviality follows by propagation.  This section states and proves
that promotion.

\begin{theorem}[b2d parameter wedge kernel triviality]
\label{thm:p12-b2d-wedge}
Let $2a < T_0 < c$, $\varepsilon := T_0 - T$, $e/2 \le R < d/2$,
$R < \sigma < \varepsilon < \varepsilon_{\max}$.  If additionally
\[
\sigma \le d/2
\qquad\text{and}\qquad
\sigma \le R + \eta,
\qquad
\eta := e - 2 \delta = \tfrac12 \log(256/243),
\]
then
$\ker L_{R, S, T_0}^{\{a, b, 2a\}} = \{0\}$.
\end{theorem}

\begin{proof}
Let $h \in L^{2}(R, S)$ with $L_{T_0} h = 0$; write $H(t) := h(T+t)$
and $l(x) := h(T - x)$.

\smallskip
\noindent\emph{Step 1 ($C_{19}$ covers $(R, \sigma)$).}
Under $\sigma \le d/2$ we have $d - \sigma \ge \sigma$; under
$\sigma \le R + \eta$ every $x \in (R, \sigma)$ satisfies
$x < \sigma \le \min\{d - \sigma, R + \eta\}$.  So
Theorem~\ref{thm:p12-b2d-core-single-slice} applies at every
$x \in (R, \sigma)$.

Because the $19 \times 19$ coefficient matrix $M_{19}$ is
non-singular, the homogeneous system forces the entire visibility
vector to vanish; in particular
\[
h(x) = 0,
\qquad l(x) = h(T - x) = 0,
\qquad h(e + x) = 0
\qquad\text{for a.e.\ } x \in (R, \sigma).
\]

\smallskip
\noindent\emph{Step 2 (kill $H$ on $(R, \sigma)$ via P1).}
For $x \in (R, \sigma)$: $d - x > d - \sigma \ge \sigma$, so
$H(d - x) = 0$.  Combined with $l(x) = 0$ from Step~1,
Theorem~\ref{thm:p12-P1} gives
\[
H(x) = 0
\qquad(x \in (R, \sigma)).
\]

\smallskip
\noindent\emph{Step 3 (homogeneous E-equation $(E_0)$ on $(R, a)$).}
For $x \in (R, \sigma)$: $H(x) = 0$ by Step~2, so the tail term
$H(x)\, \mathbf 1_{x < \sigma}$ vanishes.  For
$x \in (\sigma, a)$: $\mathbf 1_{x < \sigma} = 0$.  Hence
throughout $(R, a)$ the E-source equation reduces to the
homogeneous form
\[
p\, h(x) + r\, \mathrm{sgn}(x - d)\, h(|x - d|) - q\, h(a - x) = 0.
\tag{$E_{0}$}
\]

The reflection Steps~1--5 of Theorem~\ref{thm:p12-b2b} then apply
verbatim, using only
$e/2 \le R < e$ (satisfied since $R < d/2 < e$; elementary
$d < 2e$ i.e.\ $b - a < 2(a - d) = 2a - 2(b - a) = 4a - 2b$,
which reduces to $3b < 5a$, i.e.\ $3\log 3 < 5\log 2$,
i.e.\ $3^{3} < 2^{5}$, $27 < 32$),
and $R > \delta$ (satisfied since $R \ge e/2 > \delta$, elementary
$e > 2\delta$ iff $\log(4/3) > 2 \log(9/8)$ iff
$4/3 > (9/8)^{2}$ iff $256 > 243$).
The b2b-machinery gives
\[
h = 0
\qquad\text{a.e.\ on } (R, a).
\]
(Formally: b2b-Step~1 uses $R < e$; b2b-Steps~2--5 use $e/2 \le R$
and $R > \delta$.  b2b-Step~6, which is the only step using
$\sigma \le R$, is not invoked here because the tail is treated
separately in Step~4 below.)

\smallskip
\noindent\emph{Step 4 (kill the small tail $H$ on $(0, R)$).}
Let $0 < t < R$.  Since $R < \sigma < \varepsilon$, we have
$t < \varepsilon$, so P1 and P2 (Theorems~\ref{thm:p12-P1},
\ref{thm:p12-P2}) apply.

Compute the three $h$-values in $D(t)$:
\begin{itemize}
\item $h(t) = 0$: $t < R$ (support).
\item $h(e + t) = 0$: $e + t \in (e, e + R) \subset (R, a)$ because
  $e > R$ (from $R < d/2 < e/2$? no --- from $R < d < e + d = a$, and
  more directly $e > R$ follows from $R < d/2$ combined with
  $d/2 < e$; elementary $d < 2e$ = $27 < 32$).
  On $(R, a)$, $h = 0$ by Step~3.
\item $h(e - t) = 0$: if $e - t < R$, support; if $e - t \ge R$, then
  $e - t \in [R, e) \subset (R, a)$ where $h = 0$ by Step~3.
\end{itemize}
So $D(t) = 0$, and P2 gives
\[
H(t) - l(t) = 0.
\]

For P1: $d - t > d - R > d/2 \ge \sigma$ (since $R < d/2$ implies
$d - R > d/2$).  Hence $H(d - t) = 0$, and P1 gives
\[
H(t) + l(t) = 0.
\]

The two together force $H(t) = l(t) = 0$ on $(0, R)$.
Combined with Step~2, $H \equiv 0$ on $(0, \sigma) = (0, S - T)$:
the tail $(T, S)$ of $h$ vanishes.

\smallskip
\noindent\emph{Step 5 (reduce to b1).}
The remaining support of $h$ is in $(R, T)$.  The kernel equation
$L_{T_0} h = 0$ holds with $S$ replaced by $T$ (all deleted values
were already zero).  Theorem~\ref{thm:p12-b1} applies at
parameters $(R, T, T_0)$ under $2a < T_0 < c$ and $0 < R < T$
(satisfied since $R < d/2 < T$), giving $h = 0$.
\end{proof}

\begin{remark}[Nonemptiness of the parameter wedge]
\label{rem:p12-b2d-wedge-nonempty}
The wedge is nonempty: for any $R \in [e/2, d/2)$ and any
$\sigma \in (R, \min\{d/2, R + \eta\}]$, the hypotheses are
satisfied.  In particular, at $R = e/2$: $R + \eta = e/2 + \eta$
which is less than $d/2$ (elementary check via
$e/2 + \eta = e - \delta = \kappa < d/2$ iff $2\kappa < d$
iff $2(e - \delta) < d$ iff $2e - 2\delta < d$; using
$e = a - d$ and $\delta = 2d - a$: $2(a - d) - 2(2d - a) < d$,
$4a - 6d < d$, $4a < 7d$, $4a < 7(b - a)$, $11a < 7b$,
$11 \log 2 < 7 \log 3$, $2^{11} < 3^{7}$, $2048 < 2187$: elementary).
So at $R = e/2$, the binding constraint is $\sigma \le R + \eta$,
giving a $\sigma$-strip of width $\eta \approx 0.026$.  As $R$
increases toward $d/2$, both $R + \eta$ and $d/2$ grow, and the
$\sigma$-strip $(R, \min\{d/2, R + \eta\}]$ remains nonempty
throughout.
\end{remark}

\begin{remark}[Scope: what is and is not proven at this stage]
\label{rem:p12-b2d-wedge-scope}
Theorem~\ref{thm:p12-b2d-wedge} establishes kernel triviality on
a two-dimensional open subregion of the b2d parameter space,
strictly smaller than the full b2d chamber.  As stated here the
open frontier is
\[
\{\sigma > d/2\} \quad\text{and}\quad \{\sigma > R + \eta\}.
\]
The first is the core-both regime (where both $H(x)$ and
$H(d - x)$ are simultaneously live for $x$ near $d/2$).
The second is the post-$\eta$ visibility wall: at $x = R + \eta$,
$h(x - \eta)$ becomes live in the $Q_{18}$ source equation, so
the $19 \times 19$ closed system no longer captures all couplings
as set up.

The post-$\eta$ visibility wall is subsequently closed by the
$\eta$-wall bootstrap (\S\ref{sec:p12-b2d-eta-wall},
Theorem~\ref{thm:p12-b2d-core-single-full} and
Corollary~\ref{cor:p12-b2d-wedge-plus}), so the second constraint
$\sigma \le R + \eta$ can be removed.  At this local stage only
$\{\sigma > d/2\}$ (core-both) remained open in b2d.  The later
consolidated global result, Corollary~\ref{cor:p12-consolidated}, closes
the full b2d range because $e/2>\rho$.
\end{remark}
